\section{Stable derivators and unstable suspension}
\providecommand{\D}{\mathbb D}
\providecommand{\HoSp}{\Ho(\Sp)}

\begin{claim}{G.1}{Proved}{Every strong stable derivator has a stable cofiber localizer}
For every strong stable derivator $\D$, with the spectrum action specified below, the class
\[
 \mathcal L_{\D}=
 \{V\in\HoSp:V\otimes M=0\text{ for every }M\in\D(e)\}
\]
is a localizing subcategory and
\[
 W_{\D}=R(\mathcal L_{\D}).
\]
Consequently the reversible localizer $W_S$ of Part A is not the constant-coefficient localizer of any strong stable derivator.
\end{claim}
\begin{proof}
We use the canonical action of the derivator of spectra. Cisinski--Tabuada, \cite[Theorem A.11 and Appendix A.3]{FGHCT2011}, give the following precise universal property: for a small category $A$ and a triangulated derivator $\D$, restriction along the spectral Yoneda embedding identifies cocontinuous morphisms $\operatorname{HO}(\mathrm{Spt}_A)\to\D$ with $\D(A^{\mathrm{op}})$. For $A=e$ it assigns to $M$ the cocontinuous morphism $V\mapsto V\otimes M$ taking the sphere spectrum to $M$. The theorem is A.11 in the published paper; the opened arXiv version numbers the same statement A.15. The paper's strong triangulated derivators include the strong stable derivators considered here. No presentation of $\D$ by a model category is assumed.

This action is exact and preserves arbitrary coproducts in the spectrum variable. Its kernel for a fixed $M$ is therefore triangulated and closed under coproducts, and intersecting these kernels gives the claimed localizing subcategory. For a small category $A$, preservation of homotopy colimits identifies
\[
 (\Sigma^\infty(BA)_+)\otimes M
 \cong (p_A)_!p_A^*M,
 \qquad p_A:A\to e.
\]
Indeed the spectrum on the left is the homotopy colimit over $A$ of the constant sphere spectrum, and its action takes this diagram to the constant diagram of value $M$. In the presheaf convention one uses $A^{\mathrm{op}}$; canonical reversal of a nerve simplex gives the natural homeomorphism of classifying spaces needed to write the displayed formula with $BA$.

For $u:A\to B$, adjunction identifies the comparison on morphisms between constants with
\[
 \Hom_{\D(e)}((p_B)_!p_B^*M,N)
 \longrightarrow
 \Hom_{\D(e)}((p_A)_!p_A^*M,N).
\]
Thus the definition at I ms.~61 says, by Yoneda and for every $M$, that the specified map
\[
 \Sigma^\infty(BA)_+\otimes M
 \longrightarrow\Sigma^\infty(BB)_+\otimes M
\]
is invertible. Its cofiber is $\Sigma^\infty C(Bu)\otimes M$, since both suspension spectra and the action preserve the indicated homotopy pushout. In a stable category a map is invertible exactly when its cofiber is zero. This is the asserted equality, including empty categories and disconnected spaces. Part A proves that $W_S$ is not fixed by $RT$, whereas every $R(\mathcal L)$ is fixed by $RT$ by the Galois identities checked there; hence $W_S$ cannot be $W_{\D}$.
\end{proof}

\begin{claim}{G.2}{Proved}{Every accessible fundamental localizer has an ordinary coefficient realization}
Let $W$ be an accessible fundamental localizer, in the relative-Theorem-A convention. The derivator $\operatorname{HOT}_W$ of $W$-local spaces satisfies
\[
 W_{\operatorname{HOT}_W}=W.
\]
The realizing derivator need not be stable.
\end{claim}
\begin{proof}
Cisinski's Theorems 4.2.15 and 1.4.3 construct the corresponding accessible localization of simplicial sets; let $L:\mathcal S\to\mathcal S_W$ be its reflector. First $W$ is stable under product with any small category. For $f:A\to B$ and $C$, apply relative Theorem A over $C$ to $f\times C$. Its comma-category comparison at $c$ is $f\times(C/c)$, whose projection to $f$ is an objectwise weak equivalence because $C/c$ has a terminal object. Two out of three supplies the hypotheses. Passage through the test-category correspondence gives the same statement for products with spaces.

If $Z$ is local, then $\Map(K,Z)$ is local for every space $K$: testing a map $f\in W$ against it is the same as testing $f\times K$ against $Z$. In particular each mapping space between local objects is local. A map in $W$ consequently induces bijections on all hom-sets between constant diagrams in $\operatorname{HOT}_W$. Conversely, taking the first constant to be the terminal space gives, for every local $Z$, the bijection
\[
 [BB,Z]\longrightarrow[BA,Z].
\]
By the localization adjunction these are the maps $[LBB,Z]\to[LBA,Z]$. Yoneda in $\Ho(\mathcal S_W)$ makes the specified $LBu$ an isomorphism, which is precisely $u\in W$. This checks the previously asserted realization rather than assuming that every abstract derivator comes from localized spaces. The construction also covers the total localizer.
\end{proof}

\begin{claim}{G.3}{Proved}{Suspension for the sphere localizer is conservative, neither faithful nor full}
In the pointed $W_S$-homotopy category, suspension is conservative, neither faithful nor full. Thus it has the same behavior in any pointed derivator of $W_S$-local spaces, as in the question. These assertions also hold for the accessible smallest reversible localizer containing $\RP^\infty\to *$ constructed in Part B.
\end{claim}
\begin{proof}
For a reversible localizer, the cofiber criterion and desuspension closure give, for every pointed map $f$,
\[
 \Sigma f\in W\quad\Longleftrightarrow\quad f\in W,
\]
since $C(\Sigma f)\simeq\Sigma C(f)$. Every arrow in the formal pointed localization is represented by a roof $X\xrightarrow{f}Z\xleftarrow{w}Y$ with $w\in W$. To see the existence of such roofs, replace a map to be pushed out by its mapping-cylinder cofibration and use the homotopy pushout. Cobase change preserves $W$, because it preserves the homotopy cofiber. This moves any inverse in a finite zigzag past the following forward arrow; iteration produces the claimed roof. If its suspension is invertible, strong saturation says $\Sigma f\in W$, and the displayed implication gives $f\in W$. For $W_S$, strong saturation also follows directly from its mapping-space description established below; for the accessible localizer of Part B it follows from its localization model. Equivalently, in a local-space presentation, derived suspension is $L\Sigma$, and the same computation on a representing map proves conservativity directly.

We need an actual local target to detect failure of faithfulness. For an unpointed map $f:X\to Y$ in $W_S$, apply the pointed mapping-space functor to the cofiber sequence of $f_+:X_+\to Y_+$. Cofiber rotation gives the fibre sequence
\[
 \Map_*(\Sigma Y_+,S^n_{(2)})\longrightarrow
 \Map_*(\Sigma X_+,S^n_{(2)})\longrightarrow
 \Map_*(C(f),S^n_{(2)}).
\]
The last term is contractible, so the first arrow is an equivalence on all components. By adjunction this is the map $\Map(Y,\Omega S^n_{(2)})\to\Map(X,\Omega S^n_{(2)})$. Hence $\Omega S^n_{(2)}$ is $W_S$-local for every $n\geq2$. By the James fibre sequence checked in A.1,
\[
 S^n_{(2)}\longrightarrow\Omega S^{n+1}_{(2)}
          \longrightarrow\Omega S^{2n+1}_{(2)},\qquad n\geq2,
\]
and closure of local objects under homotopy limits, $S^n_{(2)}$ itself is local. This also proves that $W_S$ equals the class of maps inducing equivalences on mapping spaces into every $2$-local sphere: one implication has just been established, and the other follows by taking the based fibre at the constant map. G.6--G.7 supply the separate accessibility argument for this whole class.

Consider the Whitehead square $a=[\iota_2,\iota_2]:S^3\to S^2$. Its composite with $S^2\to S^2_{(2)}$ is nonzero: $\pi_3(S^2)=\ZZ$ is generated by the Hopf map $\eta$, and $a=\pm2\eta$, so localization gives the nonzero element $\pm2\in\ZZ_{(2)}$. For clarity, the coefficient two can be read directly from the attaching map of the top cell of $S^2\times S^2$: after folding the two $2$-sphere summands, the square of the degree-two cohomology generator has coefficient two. Its Hopf invariant is therefore two. The Hopf fibration identifies $\pi_3(S^2)$ with $\pi_3(S^3)=\ZZ$, and its generator has Hopf invariant one. These are the specific classical sphere facts used here; see \cite[Example 4.51, pp.~380--381, and the proof of Corollary 4J.4, p.~474]{FGHHatcher2002}.

Since $S^2_{(2)}$ is local, maps into it factor through the $W_S$-localization, so the image of $a$ there is not zero. But $\Sigma a=0$. Indeed the Whitehead product is the attaching map in the cofiber sequence $S^3\to S^2\vee S^2\to S^2\times S^2$; after suspension the inclusion of the wedge has a homotopy retraction, given by the two suspended projections and the suspension comultiplication. Exactness of this cofiber sequence makes the suspended attaching map null, and composing with the fold gives $\Sigma a=0$. Thus suspension is not faithful. For the smaller accessible reversible localizer from Part B, the same sphere remains local because its weak equivalences form a subclass of $W_S$, and the same argument applies.

Fullness already fails in the formal localization; a right adjoint to suspension need not be assumed for this conclusion. Let $i:S^2\to C(a)$ be the cofiber inclusion. Since $\Sigma a$ is nullhomotopic, the cofiber of $\Sigma a$ splits as $\Sigma S^2\vee\Sigma^2 S^3$. Under the natural identification $C(\Sigma a)\simeq\Sigma C(a)$, projection onto the first summand gives a homotopy retraction $r:\Sigma C(a)\to\Sigma S^2$ of $\Sigma i$. If localized suspension were full, there would be a morphism $g:C(a)\to S^2$ with $\Sigma g=r$. Consequently $\Sigma(gi)=\id$, and conservativity would make $gi$ invertible. But $ia=0$ in the cofiber sequence. It would follow that $a=(gi)^{-1}g(ia)=0$, contradicting the preceding paragraph. This proves nonfullness independently of the reflector construction established in G.4 and G.6--G.7.

The comparison with R9 is now precise. Its nontriviality hypothesis applies, since $S^2_{(2)}$ is a noncontractible local object. We confirm the step of GALOIS B.6 used here. In this pointed coefficient category, the identity and zero endomorphisms of $S^2_{(2)}$ are distinct. For a nontrivial acyclic group $G$, the coherent free $G$-object obtained by left Kan extension from the unique object of $BG$ has regular permutation action on the coproduct indexed by $G$. Assigning the identity to one summand and zero to the others detects a nonidentity permutation. Thus it cannot be constant, and $BG\to *$ is not a locally constant coefficient equivalence. Its suspension is weakly contractible, since $BG$ is connected and integrally acyclic. The map of double-cone covers from $BG$ to the point therefore has both cone maps and the total map in $W^{\mathrm{lc}}$, but not its intersection map, refuting Loc (6 b). The map of covers $(BG,BG,C_w,C_w)\to(BG,C_v,C_w,\Sigma BG)$ has its intersection and total maps in $W^{\mathrm{lc}}$, while $BG\to C_v$ is not, refuting Loc (6 c). These are precisely the two separate countercubes in GALOIS B.6. Hence its locally constant coefficient localizer fails both axioms even though the underlying $W_S$-homotopy suspension is conservative. The cofiber tests defining $W_S$ must not be identified with this locally constant coefficient class.
\end{proof}

\begin{claim}{G.4}{Proved}{The full sphere localizer has a small-universe reflector}
The entire $W_S$ admits an accessible reflective localization on $\mathcal S_{\mathcal U}$ in ZFC, and therefore the pointed derivator requested in Part G exists. The proof is G.6--G.7 below. Consequently every conclusion of G.3 applies to that actual derivator without an additional existence assumption.
\end{claim}
The identification of $W_S$ as coorthogonality to the set of $2$-local spheres, by itself, does not prove accessibility. G.6 separates the rational and completed-sphere conditions; G.7 constructs compact Kan models of the completed spheres and proves a countable subcomplex bound, to which B.2 applies.

\begin{claim}{G.5}{Proved}{A non-total accessible reversible localizer with full stable part}
There is a non-total accessible reversible fundamental localizer $W_*$ such that
\[
 T(W_*)=\Ho(\Sp),\qquad RT(W_*)=\operatorname{Mor}(\Cat).
\]
One may take the smallest reversible localizer containing the two maps
\[
 B\ZZ/2\longrightarrow *,\qquad B\ZZ[1/2]\longrightarrow *.
\]
In particular its intersection with finite spectra is the entire finite stable homotopy category, while $W_*$ is not fixed by $RT$.
\end{claim}
\begin{proof}
Write $\widehat S^n_2$ for the Bousfield--Kan $2$-completion of $S^n$. We first supply the completion comparison needed for the homotopy fibre sequences
\begin{equation}\label{eq:G-completed-EHP}
 \widehat S^n_2\longrightarrow\Omega\widehat S^{n+1}_2
       \longrightarrow\Omega\widehat S^{2n+1}_2,
 \qquad n\geq2.
\end{equation}
Let $F$ be the homotopy fibre of the integral James--Hopf map
$H:\Omega S^{n+1}\to\Omega S^{2n+1}$. The usual nullhomotopy of $HE$ gives a map $S^n\to F$. The precise $2$-local EHP theorem checked in A.1 says that this map becomes an equivalence after localization at $2$, hence is a mod-$2$ homology equivalence. All three spaces $F$, $\Omega S^{n+1}$ and $\Omega S^{2n+1}$ are simply connected and have finitely generated homotopy groups: the loop-space assertions follow from Serre finiteness, and those for $F$ follow from the fibre long exact sequence (for simple connectedness also use $\pi_2\Omega S^{2n+1}=0$). These are exactly the finite nilpotence hypotheses in \cite[Theorem 4.5.2 and Theorem 11.2.6]{FGHMayPonto2012}. The latter theorem, including its final simply connected-base clause on page~220, therefore identifies the completion of $F$ with the fibre of the completed map $H$. Proposition~11.2.4 of the same source identifies the completions of the connected loop spaces with the loop spaces of the completed spheres. Finally $S^n\to F$ induces an equivalence on completions by Theorems~11.1.1--11.1.2, since it is a mod-$2$ homology equivalence. This proves \eqref{eq:G-completed-EHP} with the specified maps. Equivalently, their homotopy-group comparison can be checked by the natural formula $\pi_i(X^{\wedge}_2)=\pi_i(X)\otimes\ZZ_2$ and exactness of tensoring with the torsion-free group $\ZZ_2$; Theorem~11.1.2 and Lemma~10.1.4 justify that formula under precisely the finiteness conditions just verified.

Define
\[
 \widehat{\mathcal N}_2
 =\{X\in\Ho(\mathcal S_*):\Map_*(X,\widehat S^n_2)\simeq *
                       \text{ for every }n\geq2\}.
\]
It is closed under pointed homotopy colimits. Sequence \eqref{eq:G-completed-EHP} shows that it is closed under desuspension: if $\Sigma X$ is in the class, both mapping spaces into the loop-space terms are contractible, and their homotopy fibre is $\Map_*(X,\widehat S^n_2)$. Suspension closure follows directly by adjunction. Cofiber rotation then gives two out of three for cofiber sequences. Thus the checked construction A.1--A.3 applies and gives a reversible fundamental localizer
\[
 \widehat W_2=\{f:C(f)\in\widehat{\mathcal N}_2\}.
\]
The completed spheres are actual $\widehat W_2$-local objects: repeat the cofiber-rotation argument in G.3 to obtain locality of their loop spaces, and then use \eqref{eq:G-completed-EHP}. In particular this is an all-component mapping-space assertion, not merely a test on the fibre over the constant map.

The completion map $S^n\to\widehat S^n_2$ is a mod-$2$ homology equivalence, so $\widehat S^n_2$ is nilpotent with mod-$2$ cohomology in only degrees $0,n$. Miller's Theorem~C, in the exact form verified in A.4, puts $B\ZZ/2$ in $\widehat{\mathcal N}_2$. Also $B\ZZ[1/2]$ is the mapping telescope of the degree-two self maps of $S^1$. Its reduced mod-$2$ homology vanishes, since in degree one the direct system is $\FF_2\xrightarrow{0}\FF_2\xrightarrow{0}\cdots$ and all higher groups vanish. Each $\widehat S^n_2$ is mod-$2$ homology-local by \cite[Theorems 11.1.1--11.1.2]{FGHMayPonto2012}; hence $\Map_*(B\ZZ[1/2],\widehat S^n_2)$ is contractible. Both generators therefore belong to $\widehat W_2$.

Let $W_*$ be their smallest reversible closure. Part B proves that it is accessible, and minimality gives $W_*\subseteq\widehat W_2$. Obstruction theory on the two-dimensional complex $\RP^2$ gives
\[
 [\RP^2,\widehat S^2_2]_*\cong
 H^2(\RP^2;\ZZ_2)\cong\ZZ/2,
\]
using $\pi_1\widehat S^2_2=0$ and $\pi_2\widehat S^2_2=\ZZ_2$. Thus $\RP^2$ is not in $\widehat{\mathcal N}_2$, and $\RP^2\to *$ is not in $\widehat W_2$ by suspension and desuspension closure. This proves that $W_*$ is non-total.

Its stable part nevertheless contains both $\mathbb S/2$ and $\mathbb S[1/2]$. The first follows from $B\ZZ/2\to *$ by the independently checked Kahn--Priddy argument of A.5. For the second, suspension spectra preserve the circle telescope and give
\[
 \Sigma^\infty B\ZZ[1/2]\simeq\Sigma\mathbb S[1/2].
\]
The cofibres of the two generating maps have one further suspension, which is immaterial in a localizing subcategory. To see directly that these two spectra generate every spectrum, the cofiber of $\mathbb S\to\mathbb S[1/2]$ is the homotopy colimit of the cofibres of $2^r:\mathbb S\to\mathbb S$. Each such cofiber $\mathbb S/2^r$ belongs to $\operatorname{loc}(\mathbb S/2)$, by induction using the triangles for successive powers of two. The cofiber triangle therefore puts $\mathbb S$ in $\operatorname{loc}(\mathbb S/2,\mathbb S[1/2])$, and the sphere generates $\Ho(\Sp)$. Hence $T(W_*)=\Ho(\Sp)$ and $RT(W_*)$ is the total localizer, whereas $W_*$ is non-total as proved above.
\end{proof}

\begin{claim}{G.6}{Proved}{The sphere localizer is the intersection of completed-sphere and rational equivalences}
Every rational Eilenberg--Mac Lane space $K(\QQ,d)$, $d\geq0$, is $W_S$-local. Consequently, with $\widehat{\mathcal N}_2$ and $\widehat W_2$ as in G.5,
\[
 \Nul=\widehat{\mathcal N}_2\cap
       \{X:\widetilde H_*(X;\QQ)=0\},
 \qquad
 W_S=\widehat W_2\cap W_{\QQ},
\]
where $W_{\QQ}$ denotes rational homology equivalences.
\end{claim}
\begin{proof}
Fix an odd integer $n\geq3$. By G.3, the loop space $V_n=\Omega S^n_{(2)}$ is $W_S$-local. Its multiplication gives a pointed squaring map $s:V_n\to V_n$, $s(v)=v\cdot v$. On every homotopy group $\pi_i(V_n)$, $i\geq1$, this map induces multiplication by two: the group operation induced by the loop-space multiplication agrees with the usual operation on homotopy groups. We use this squaring map on the loop space, so no assertion about the action of a degree-two sphere map in unstable homotopy groups is required.

Consider the homotopy inverse limit
\[
 P_n=\holim\bigl(V_n\xleftarrow{s}V_n\xleftarrow{s}V_n\xleftarrow{s}\cdots\bigr).
\]
It is $W_S$-local, since all its vertices are local. The homotopy groups of a finite simply connected complex are finitely generated, and the rationalization of an odd sphere is $K(\QQ,n)$; see \cite[Theorem 4.5.2 and Corollary 6.7.3]{FGHMayPonto2012}. Localization of homotopy groups therefore gives
\[
 \pi_{n-1}(V_n)=\ZZ_{(2)},
 \qquad
 \pi_i(V_n)\text{ is a finite }2\text{-group for }i\ne n-1.
\]
This includes zero groups in the lower dimensions. The inverse system given by multiplication by two on a finite $2$-group is pro-zero, so both its limit and its first derived limit vanish. For $A=\ZZ_{(2)}$, the limit of the tower under multiplication by two is zero and
\begin{equation}\label{eq:G-limone-rational}
 \lim{}^1(A\xleftarrow{2}A\xleftarrow{2}\cdots)
       \cong \ZZ_2/\ZZ_{(2)}.
\end{equation}
Here is the full algebraic calculation. The derived limit is the cokernel of
\[
 D:\prod_{i\geq0}A\longrightarrow\prod_{i\geq0}A,
 \qquad D((a_i))=(a_i-2a_{i+1}).
\]
The convergent $2$-adic series
$\Phi((b_i))=\sum_{i\geq0}2^i b_i$
induces a map from this cokernel to $\ZZ_2/A$, since
$\Phi(D(a))=a_0$. It is onto by the binary expansion of a $2$-adic integer. If $\Phi(b)=a_0\in A$, define
\[
 a_j=2^{-j}\left(a_0-\sum_{i<j}2^i b_i\right).
\]
The quantity in parentheses has $2$-adic valuation at least $j$, so $a_j\in\QQ\cap\ZZ_2=A$, and $a_j-2a_{j+1}=b_j$. This proves injectivity and hence \eqref{eq:G-limone-rational}. The ordinary limit vanishes because $\bigcap_j2^j A=0$.

The homotopy inverse-limit exact sequence now shows
\[
 P_n\simeq K(D_2,n-2),\qquad D_2=\ZZ_2/\ZZ_{(2)}.
\]
There is no extra component: $V_n$ is connected and $\pi_1(V_n)=0$ for $n\geq3$, so the component term $\lim^1\pi_1(V_n)$ is zero. In the case $n=3$ the single nonzero group occurs in degree one, and the same exact sequence identifies it with the abelian group $D_2$. In all larger cases the displayed calculation identifies all homotopy groups; a CW approximation gives the asserted Eilenberg--Mac Lane homotopy type.

The group $D_2$ is a nonzero rational vector space. Multiplication by an odd integer is invertible on both groups in the quotient. Multiplication by two on $D_2$ is onto because one may subtract an integer from a $2$-adic integer to make it divisible by two; it is one-to-one because $2a\in\ZZ_{(2)}$ and $a\in\ZZ_2$ imply $a\in\QQ\cap\ZZ_2=\ZZ_{(2)}$. Finally $D_2\ne0$, since $\ZZ_2$ is uncountable and $\ZZ_{(2)}$ is countable. Choose a one-dimensional rational direct summand. The induced maps make $K(\QQ,n-2)$ a homotopy retract of $P_n$, so it is local. Odd $n$ can be arbitrarily large. Taking loop spaces, which are homotopy limits of local objects, gives locality of $K(\QQ,d)$ for every $d\geq0$.

It follows that every map in $W_S$ is a rational homology equivalence. Indeed locality of all these Eilenberg--Mac Lane spaces gives isomorphisms on ordinary rational cohomology in all nonnegative degrees. The universal coefficient theorem over the field $\QQ$ identifies these maps with the linear duals of the corresponding homology maps; a linear map of vector spaces is an isomorphism if its dual is, including for infinite-dimensional vector spaces. Thus $W_S\subseteq W_{\QQ}$. If $X\in\Nul$, A.3 makes $X\to *$ a $W_S$-equivalence, and so $X$ is rationally acyclic.

For every $n\geq2$ the arithmetic fracture square of the nilpotent $2$-local sphere is
\[
 S^n_{(2)}\simeq
 (S^n_{(2)})_{\QQ}\mathop{\times}\limits^h_{(\widehat S^n_2)_{\QQ}}
                    \widehat S^n_2;
\]
see \cite[Theorem 13.1.4]{FGHMayPonto2012}. If $X$ is rationally acyclic, mapping it into either rational vertex gives a contractible space. Hence this square gives
\[
 \Map_*(X,S^n_{(2)})\simeq\Map_*(X,\widehat S^n_2).
\]
The preceding rational-acyclicity conclusion now proves both inclusions in the asserted equality of null classes. Applying that equality to $C(f)$ proves the equality of localizers, since $\widetilde H_*(C(f);\QQ)=0$ is exactly the rational homology equivalence condition for $f$. This last statement includes the empty-source cases through the cofiber convention $C(f)=\operatorname{cofib}(f_+)$ used in A.3.
\end{proof}

\begin{claim}{G.7}{Proved}{The full sphere localizer is accessible in ZFC}
The null classes of G.5 and Part A satisfy
\[
 \widehat{\mathcal N}_2
       =\widehat{\mathcal N}_2^{\aleph_1},
 \qquad \Nul=\Nul^{\aleph_1}.
\]
Consequently both $\widehat W_2$ and $W_S$ are accessible reversible fundamental localizers in ZFC. A generating set for each is given by the category-of-simplices models of the maps between countable CW complexes whose cofibers lie in its displayed null class. In particular $W_S$ itself is the accessible counterexample asked for in Part B, and the pointed localization derivator in Part G exists in the stipulated small universe.
\end{claim}
\begin{proof}
We give a strict compact model and a compactness argument for small subcomplexes. The proof does not identify the homotopy category of an arrow category with the arrow category of a homotopy category.

\emph{Compact Kan models.}
Let $K=\Delta^n/\partial\Delta^n$, $n\geq2$, and let $R$ be the reduced free simplicial $\FF_2$-module functor, regarded as a monad on pointed simplicial sets. The standard Bousfield--Kan resolution has cosimplicial terms
\[
 X^q=R^{q+1}K\qquad(q\geq0).
\]
The precise construction and its identification with $\FF_2$-completion are in \cite[\S7.7, pp.~1032--1033]{FGHBousfield2003}. Each simplicial degree of each $X^q$ is finite, since $K$ has finitely many simplices in each degree and a free vector space on a finite set over $\FF_2$ is finite. Also each $X^q$ is a Kan complex, being a simplicial abelian group.

Here is a direct check that this cosimplicial resolution is Reedy fibrant. Its codegeneracies $s^i$ are linear for the outer $R$-module structures, as are all cofaces $d^i$ with $i\geq1$. A matching tuple in degree $q$ is determined by compatible elements $(a_0,\ldots,a_{q-1})$ in degree $q-1$. Start with $c=0$. If $s^ic=a_i$ for $i<j$, replace $c$ by
\[
 c+d^{j+1}(a_j-s^jc).
\]
The identities $s^jd^{j+1}=\id$ and $s^id^{j+1}=d^js^i$ for $i<j$ show that this corrects the $j$th equation and preserves the previous ones: the latter assertion uses $s^i a_j=s^{j-1}a_i$. This construction is linear and simplicial. Thus each matching map is a split surjection of simplicial abelian groups, hence a Kan fibration. In cosimplicial degree zero the assertion is just the Kan property already noted.

Totalization is right Quillen for the Reedy structure; see the statement immediately before \cite[Proposition 8.1, p.~1034]{FGHBousfield2003}, or check the horn lifting adjunction against the Reedy cofibrations $\Delta^\bullet\otimes\Lambda^k[m]\to\Delta^\bullet\otimes\Delta[m]$. Therefore
\[
 Z_n=\operatorname{Tot}(R^{\bullet+1}K)
\]
is Kan and is a model of $\widehat S^n_2$. In simplicial degree $m$, its set of simplices is the closed compatibility subspace of
\[
 \prod_{q\geq0}
 \operatorname{Hom}_{\mathrm{sSet}}
     (\Delta^q\times\Delta^m,R^{q+1}K).
\]
Each factor is finite: the product of the two standard simplices has finitely many nondegenerate simplices. Give each factor the discrete topology. The displayed product is compact metrizable, its compatibility equations are closed, and all face and degeneracy operators are continuous. Thus $Z_n$ is a simplicial compact metrizable Hausdorff space with underlying Kan complex. This argument supplies an actual strict simplicial model, not just compact topologies on homotopy classes.

\emph{A bounded extension lemma.}
Suppose $Z$ is any pointed Kan complex whose levels have compact metrizable Hausdorff topologies and continuous simplicial operators. For a pointed simplicial set $A$, the set
\[
 H_Z(A)=\operatorname{Hom}_{\mathrm{sSet}_*}(A,Z)
\]
is a closed subspace of the product of the target levels indexed by simplices of $A$, hence is compact Hausdorff. For countable $A$ it is metrizable. Restrictions are continuous.

Fix a pointed simplicial set $X$ and a countable pointed simplicial subset $A\subseteq X$. Let $C_A$ be the image of $H_Z(X)\to H_Z(A)$. This is compact and closed. For $A\subseteq L\subseteq X$ obtained by adjoining finitely many nondegenerate simplices and all their faces, let $C_{L,A}$ be the image of $H_Z(L)\to H_Z(A)$. Then
\begin{equation}\label{eq:G-compact-image-intersection}
 C_A=\bigcap_L C_{L,A}.
\end{equation}
Indeed a compatible assignment on $A$ belongs to every image on the right precisely when every finite set of additional simplex equations can be solved. In the compact product of all the target coordinates indexed by $X$, those equations and the prescribed coordinates on $A$ are closed conditions with the finite intersection property. Compactness gives a simultaneous solution, which is the required map $X\to Z$.

Choose a countable basis of $H_Z(A)$. The basic open sets whose closures miss $C_A$ cover its complement: compact metrizable spaces are regular, so the basis may be chosen to have this property. For each such open $U$, equation \eqref{eq:G-compact-image-intersection} and compactness of $\overline U$ show that finitely many of the sets $C_{L,A}$ already have intersection disjoint from $\overline U$. Their finite union of added simplices gives one countable enlargement $L_U$ with this property. The union of these $L_U$ is a countable enlargement $A'\subseteq X$ such that
\begin{equation}\label{eq:G-one-step-extension}
 \operatorname{im}(H_Z(A')\to H_Z(A))=C_A.
\end{equation}
The same construction can be performed simultaneously for a countable family of such targets by taking the countable union of the enlargements.

Apply this to the countable family
\[
 Z_{n,m}=\operatorname{Map}_*(\Delta^m/\partial\Delta^m,Z_n),
       \qquad n\geq2,\quad m\geq0,
\]
where $m=0$ means $Z_{n,0}=Z_n$. These are Kan models of $\Omega^m Z_n$ and again have compact metrizable levels with continuous operators: each level is a closed subset of a finite product of target levels, because the simplicial sphere smashed with a standard simplex is finite. Starting at any countable $A_0\subseteq X$, iterate \eqref{eq:G-one-step-extension} to obtain countable
\[
 A_0\subseteq A_1\subseteq A_2\subseteq\cdots,
 \qquad A_\infty=\bigcup_i A_i.
\]
Every map $h:A_\infty\to Z_{n,m}$ extends to $X$. To see this, its restriction to $A_i$ extends globally because it is the restriction of its map on $A_{i+1}$. The sets of global extensions of $h|_{A_i}$ are nonempty nested closed subsets of the compact space $H_{Z_{n,m}}(X)$. Their intersection supplies an extension of $h$ itself.

Now assume $X\in\widehat{\mathcal N}_2$. Each such global extension is pointed nullhomotopic, since $\Map_*(X,Z_{n,m})$ is weakly contractible. Its restriction makes $h$ nullhomotopic. Thus
\[
 [A_\infty,\Omega^m Z_n]_* = *
       \quad\text{for all }n\geq2,\ m\geq0.
\]
The $m=0$ assertion makes each $\Map_*(A_\infty,Z_n)$ connected; the other assertions make every homotopy group at its basepoint zero. These mapping spaces are Kan, so they are weakly contractible. Hence $A_\infty\in\widehat{\mathcal N}_2$.

We have proved that every countable simplicial subset of a member of $\widehat{\mathcal N}_2$ is contained in a countable member of that class. Realization converts this to the CW version of B.1: the realization of a countable simplicial set has countably many cells, and every countable collection of cells in a realization is contained in the realization of a countable simplicial subset. Starting with a simplicial model of an arbitrary CW homotopy type and then using the model-invariance argument of B.1 gives
$\widehat{\mathcal N}_2=\widehat{\mathcal N}_2^{\aleph_1}$. This step does not assume that the singular complex of a countable CW complex is countable.

\emph{The rational condition and conclusion.}
The class of connected rationally acyclic spaces has the same countable enlargement property. Given a countable subcomplex of a rationally acyclic $X$, there are only countably many finite rational chains and cycles on it. Adjoin the finite supports of bounding chains in $X$ for all its reduced cycles, and repeat countably. Every cycle in the resulting union occurs at some stage and bounds at the next; the union is therefore rationally acyclic. This includes reduced degree zero and so also ensures connectedness.

If $X\in\Nul$, G.6 says that $X$ is both completed-sphere-null and rationally acyclic. Alternate the two countable enlargement constructions just proved. Their common sequential union belongs to both classes, because both are closed under pointed homotopy colimits. It therefore belongs to $\Nul$, again by G.6. Thus $\Nul=\Nul^{\aleph_1}$. B.2 applies to both null classes, proving the stated accessibility and the explicit generating sets. Cisinski's localization construction used in G.2 now gives the reflector and its pointed derivator in $\mathcal U$, so the suspension and coefficient conclusions of G.3 are unconditional ZFC statements for the full $W_S$.
\end{proof}

